Returning to the context of 5.6, the algebraic theory in $\mathcal{A}$ giving rise to $V: \mathcal{C} \longrightarrow \mathcal{A}$ is given by $(X, s)\hat{T} = (XT, \hat{s})$ where $\hat{s}$ is the cartesian lift of the family

$$\{KT \xrightarrow{g} (L, t) \mid \text{there exists } \xi \text{ with } (L, t, \xi) \in \mathcal{C} \text{ and } g: (KT, K\mu) \longrightarrow (L, \xi) \text{ a } \mathbf{T}\text{-homomorphism}\}.$$

$K\eta: (K, s) \longrightarrow (K, s)\hat{T}$ is admissible, and if $\alpha: (K_1, s_1) \longrightarrow (K_2, s_2)\hat{T}$ and $\beta: (K_2, s_2) \longrightarrow (K_3, s_3)\hat{T}$ are admissible, then so is $\alpha \circ \beta: (K_1, s_1) \longrightarrow (K_3, s_3)\hat{T}$, thereby providing the $\eta$ and $\circ$ for the algebraic theory $\hat{T} = (\hat{T}, \eta, \circ)$. The proof follows from applying Proposition 5.10.
